\section{Introduction}

Technical standards can increase interoperability and enlarge the set of products that can use a common technological improvement. But firms do not passively accept the competitive environment created by a standard. When a broader standard makes an improvement in the common layer more useful to rival products, a firm may respond by moving scarce innovative effort toward technologies that remain proprietary. The policy question is therefore not only how much interoperability a standard creates directly, but also how standard scope changes the composition of innovation.

We study that feedback in a deliberately stripped-down industrial-organization model. A regulator first chooses the scope of a mandatory common technical standard. Two firms then allocate a fixed R\&D capacity between common-layer and proprietary innovation, after which they compete in differentiated Bertrand prices. Common-layer R\&D improves the innovating firm's product and, to an extent determined by the mandated scope, the rival's standard-compliant product. Proprietary R\&D benefits only the innovating firm. Scope therefore changes the relative appropriability of two uses of a fixed innovation budget without directly changing product substitutability.

The first result is an endogenous portfolio response. In the quadratic baseline, broader scope strictly reduces private common-layer R\&D and increases proprietary R\&D. The directional result is more general but weaker: for any differentiable increasing strictly concave innovation technology, the private common-layer allocation is globally nonincreasing in scope, while a coordinated symmetric allocation is globally nondecreasing. If two compared scope choices both induce interior optima, the corresponding order inequalities are strict. We do not claim a pointwise derivative sign for arbitrary general technologies. This distinction matters because corner solutions can generate flat intervals and differentiability of the innovation technology alone does not justify differentiating the optimizer with respect to policy.

The main policy result is deliberately baseline-specific. If the R\&D composition is held fixed at any positive symmetric common-layer allocation, welfare is increasing in scope and complete standardization is optimal. Complete standardization is also optimal when the regulator chooses scope together with the symmetric R\&D composition while differentiated Bertrand pricing remains decentralized. That benchmark is constrained and is not an unconstrained first best. With decentralized endogenous R\&D composition, broader scope induces firms to retreat from common innovation. In the quadratic baseline, sufficiently strong product-market rivalry makes the reduced policy objective strictly concave with a unique interior optimum. Figure~\ref{fig:regime-map} summarizes the exact rivalry threshold. Selective standardization is therefore a scope-only second-best response to an innovation-composition distortion, not a claim that interoperability is intrinsically socially costly.

The contribution is deliberately result-level rather than thematic. The closest technical-standards comparison is \citet{Llanes2024}, who studies incentives to develop technologies for inclusion in a standard together with patent pools, price caps, and cooperative R\&D. Related work shows that compatibility and standard-setting institutions affect innovation incentives \citep{BondSmith2019,MaruyamaZennyo2017,HeywoodWangYe2022}, while endogenous product design can overturn compatibility conclusions \citep{Boom2001} and government standardization can interact with product characteristics \citep{Ruiz2004}. Our policy object is different: a regulator chooses a continuous standardization scope that changes cross-product transferability before two firms reallocate a fixed common-versus-proprietary R\&D portfolio. \citet{AcemogluGanciaZilibotti2012} already characterize constrained optimal standardization in a dynamic growth environment, so a generic second-best case for limiting standardization is not novel here. Likewise, \citet{BryanLemus2017} develop a theory of the direction of innovation with scarce research resources and underappropriation; endogenous innovation direction itself is therefore not our novelty. The paper's contribution is the combination of regulator-chosen scope, endogenous common-versus-proprietary portfolio reallocation, differentiated-product rivalry, and the exact complete-to-selective policy reversal in the quadratic baseline.

The mechanism is intended for settings with a meaningful common interface or standardized technological layer and a distinct proprietary layer. Digital protocols and APIs and modular industrial interfaces are natural interpretations. Evidence from modular product development is consistent with innovative attention shifting toward internal interfaces and modules when external interfaces are standardized \citep{ChenLiu2005}. Recent firm-level evidence also shows that firms technologically closer to a newly adopted standard increase R\&D more strongly in competitive markets and that standardization can create a common technological basis that facilitates later catch-up \citep{BergeaudSchmidtZago2026}. We use these findings only to motivate the interaction among standardization, competition, and innovation; they are not causal validation of the portfolio mechanism modeled here.

The rest of the paper proceeds as follows. Section~2 defines the game and the all-history price-continuation restriction. Section~3 characterizes private R\&D allocation and derives the selective-standardization threshold. Section~4 derives welfare and compares the decentralized policy problem with fixed-symmetric and coordinated-symmetric benchmarks without claiming an unconstrained first best. Section~5 states only the robustness results authorized by the theory freeze. Section~6 positions the contribution relative to compatibility, standards-and-innovation, innovation-direction, and R\&D-allocation literatures. The Appendix provides global price-continuation, corner/KKT, threshold, and verification details.
